\chapter{Experiments and Results}
\label{sec:experiments_and_results}

This chapter presents the experimental setup and the quantitative results from the implemented pipeline. The focus is on reporting what was tested and what the models achieved. Deeper interpretation of why certain components helped or failed is kept for Chapter~6.


\section{Experimental Protocol}
\label{sec:experimental_protocol}

All experiments were evaluated at patient level. For each patient \(i\), the model produced one final prediction \(\hat{y}_i\), which was compared with the corresponding patient-level target \(y_i\). Slice-level scores were only used as intermediate outputs inside the model and were not evaluated separately, since no slice-level target labels were available.


\subsection{Internal Five-Fold Cross-Validation}
\label{sec:internal_five_fold_cv_results}

The internal mucus-score experiments used the five-fold cross-validation setup described in Chapter~4. In each fold, the model was trained on the training patients and evaluated on the held-out validation patients. The same patient-level fold assignments were used across the main internal experiments so that different model configurations could be compared under the same validation structure. Results from the internal experiments are reported as mean \(\pm\) standard deviation across the five folds.

Baseline experiments were trained using the raw mucus-score target, while others used the log-transformed target described in Chapter~4. For easier interpretation, the reported error metrics are transformed back to the raw mucus-score scale whenever possible, so that MAE and RMSE can be read in terms of the original burden score.

The internal experiments compare slice handling, input representation, feature fusion, and backbone choice. The two-window ResNet18 model is used as the main reference setting. Additional experiments evaluate the representative slice-selection strategy, three-channel input representations, ORB-BoVW feature fusion, and DenseNet121 as a comparison backbone.


\subsection{Transfer-Learning Evaluation}
\label{sec:transfer_learning_evaluation_protocol}

Transfer-learning experiments are reported separately from the internal five-fold experiments. These experiments evaluate how models behave when trained or initialized across the mucus dataset and the external MosMed dataset. For these experiments, the tables specify which dataset is used as the source domain and which dataset is used as the target domain.

The MosMed experiments are transfer-learning and domain-shift experiments, not external validation of mucus-burden prediction. MosMed uses a different case-level CT severity target, so mucus-to-MosMed evaluates adaptation to the MosMed severity task, while MosMed-to-mucus evaluates adaptation back to the internal mucus-burden task.

The transfer results are reported as average performance across the five folds. The target was modelled on the log-transformed scale, and predictions were transformed back to the raw target scale for reporting. This makes the reported MAE and RMSE easier to interpret in terms of the original score values.

In the transfer tables, zero-shot evaluation refers to applying a source-trained model to the target domain without target-domain training. Fine-tuned evaluation refers to adapting the source-initialized model using target-domain training data before evaluation.



\section{Evaluation Metrics}
\label{sec:results_metrics}

The models are evaluated using patient-level regression metrics. Let \(y_i\) denote the observed target for patient \(i\), and let \(\hat{y}_i\) denote the corresponding patient-level prediction. For a validation set with \(N\) patients, the mean absolute error is defined as
\[
\mathrm{MAE}
=
\frac{1}{N}
\sum_{i=1}^{N}
|y_i-\hat{y}_i|.
\]
MAE measures the average absolute size of the prediction error. Lower values indicate better performance.

The root mean squared error is defined as
\[
\mathrm{RMSE}
=
\sqrt{
\frac{1}{N}
\sum_{i=1}^{N}
(y_i-\hat{y}_i)^2
}.
\]
RMSE also measures prediction error, but it gives larger errors more weight than MAE because the residuals are squared before averaging. Lower RMSE is therefore better. MAE, RMSE, and \(R^2\) are commonly used regression evaluation metrics, but they emphasize different aspects of prediction error and should therefore be interpreted together rather than in isolation \cite{chicco2021coefficient}.

The coefficient of determination is computed as
\[
R^2
=
1 -
\frac{
\sum_{i=1}^{N}(y_i-\hat{y}_i)^2
}{
\sum_{i=1}^{N}(y_i-\bar{y})^2
},
\]
where \(\bar{y}\) is the mean target value in the evaluated set. Higher \(R^2\) is better, with \(1\) representing perfect prediction. However, \(R^2\) should be interpreted relative to the mean-prediction reference. An \(R^2\) value close to zero means that the model explains little additional variation beyond predicting the mean target value, while a negative \(R^2\) means that the model performs worse than this reference for that evaluation set. Since the validation folds are small in this thesis, \(R^2\) is interpreted together with MAE, RMSE, and Spearman correlation rather than as the only summary of performance \cite{chicco2021coefficient}.

Spearman correlation is used to measure rank association between the observed targets and the model predictions \cite{spearman1904proof}. This metric is useful because it shows whether the model tends to order patients correctly by predicted burden, even when the absolute prediction errors are still large. Higher Spearman correlation indicates stronger agreement in ranking.

In addition to reporting mean \(\pm\) standard deviation across the five folds, an approximate confidence interval is reported for the proposed model. For a metric value \(m_f\) computed on fold \(f\), with fold mean \(\bar{m}\), sample standard deviation \(s_m\), and \(F=5\) folds, the approximate \(95\%\) confidence interval is computed using a Student \(t\)-interval:
\[
\bar{m}
\pm
t_{0.975,F-1}
\frac{s_m}{\sqrt{F}}.
\]
Here, \(t_{0.975,F-1}\) is the \(97.5\)th percentile of the \(t\)-distribution with \(F-1\) degrees of freedom. This interval gives a compact indication of fold-level uncertainty for the proposed model. Because it is based on only five folds, it should be interpreted cautiously and used as a robustness summary rather than as a definitive population-level uncertainty estimate \cite{openstax_t_interval}.


\section{Dataset and Fold Summary}
\label{sec:dataset_fold_summary}

The internal experiments were performed on \(N=32\) included patient cases. Each case had one CT volume and one patient-level mucus-score target. The raw target values ranged from \(0.0\) to \(7.5\), with a mean of \(2.688\) and a standard deviation of \(2.313\). The median raw target was \(2.0\), with an interquartile range from \(1.0\) to \(4.0\). The log-transformed target ranged from \(0.000\) to \(0.875\), with a mean of \(0.361\) and a standard deviation of \(0.303\).

The fold sizes were slightly uneven because of the small cohort size. Folds~0 and~1 contained seven patients each, while folds~2, 3, and~4 contained six patients each. Table~\ref{tab:dataset_fold_summary} summarizes the raw target distribution across folds.

\begin{table}[H]
\centering
\caption{Patient counts and raw target distribution across the five internal folds.}
\label{tab:dataset_fold_summary}
\small
\begin{tabular}{cccccc}
\toprule
Fold & Patients & Min & Max & Mean & Median \\
\midrule
0 & 7 & 0.0 & 5.0 & 2.571 & 2.00 \\
1 & 7 & 0.0 & 5.5 & 2.286 & 2.00 \\
2 & 6 & 0.0 & 7.5 & 2.417 & 2.00 \\
3 & 6 & 0.0 & 7.0 & 2.917 & 2.25 \\
4 & 6 & 0.0 & 7.0 & 3.333 & 2.25 \\
\bottomrule
\end{tabular}
\end{table}

The number of slices per CT volume varied from \(47\) to \(716\), with a median of \(241\) slices and an interquartile range from \(151.25\) to \(337\). This variation is one reason why the pipeline uses a slice-selection strategy before model training.


\section{Internal Experiments}
\label{sec:internal_experiments}


\subsection{Baseline Experiments}
\label{sec:baseline_experiments}
The baseline experiments use simpler reference configurations than the main pipeline. They are included to show how the final method compares with more basic slice-handling and input settings. The baseline experiments compare two slice-handling settings using a ResNet18 model with two-window CT input. The uniform-sampling baseline samples slices evenly from the CT volume, while the other ResNet18 model with the same two-window CT input uses the representative slice-selection strategy described in Chapter~4. Because this baseline does not use the full final pipeline configuration, it should be interpreted as a simpler reference setting rather than as a direct ablation of only one component.



\begin{table}[H]
\centering
\caption{Two-window baseline comparison between uniform slice sampling and the representative slice-selection strategy. Values are reported as mean \(\pm\) standard deviation across patient-level folds.}
\label{tab:baseline_results}
\small
\begin{tabular}{lcc}
\toprule
Metric & Uniform sampling & Slice sampling strategy \\
\midrule
MAE & \(1.922 \pm 0.326\) & \(1.760 \pm 0.154\) \\
RMSE & \(2.268 \pm 0.301\) & \(2.223 \pm 0.209\) \\
\(R^2\) & \(-0.009 \pm 0.064\) & \(0.003 \pm 0.109\) \\
Spearman & \(0.202 \pm 0.309\) & \(0.306 \pm 0.330\) \\
\bottomrule
\end{tabular}
\end{table}

The representative slice-selection strategy performed better than uniform sampling on MAE and Spearman correlation, while RMSE and \(R^2\) changed only modestly. This suggests that the selected slices may have provided a more useful input set for patient-level prediction, although the small fold sizes mean that the difference should be interpreted cautiously.


\subsection{Input Representation and Feature Fusion}
\label{sec:input_representation_feature_fusion_results}

Table~\ref{tab:internal_representation_fusion_results} reports the main internal comparison between the two-window input, the three-channel input, and the ORB-BoVW fusion model. These experiments use the main pipeline configuration described in Chapter~4, including the representative slice-selection settings, patient-level aggregation setup, and training configuration summarized in the methodology tables. This makes the rows in this table more directly comparable to each other than to the simpler baseline settings in Table~\ref{tab:baseline_results}.

The models in this comparison were trained using the log-transformed target. For reporting, predictions were transformed back to the raw mucus-score scale so that MAE and RMSE are easier to interpret.

\begin{table}[H]
\centering
\caption{Main internal five-fold results for input representation and feature fusion. Models used the main configuration from Chapter~4, were trained on the log-transformed target, and were transformed back to the raw mucus-score scale for reporting. Values are mean \(\pm\) standard deviation across patient-level folds.}
\label{tab:internal_representation_fusion_results}
\small
\setlength{\tabcolsep}{4pt}
\begin{tabularx}{\textwidth}{@{}>{\raggedright\arraybackslash}Xcccc@{}}
\toprule
Experiment & MAE & RMSE & \(R^2\) & Spearman \\
\midrule
ResNet18, two-window input & \(1.451 \pm 0.395\) & \(1.858 \pm 0.434\) & \(0.126 \pm 0.161\) & \(0.363 \pm 0.239\) \\
ResNet18, three-channel input & \(1.386 \pm 0.456\) & \(1.826 \pm 0.394\) & \(0.186 \pm 0.394\) & \(0.387 \pm 0.256\) \\
\textbf{ResNet18-3CH+BoVW (Proposed model)} & \(\mathbf{1.244 \pm 0.297}\) & \(\mathbf{1.712 \pm 0.505}\) & \(\mathbf{0.236 \pm 0.295}\) & \(\mathbf{0.498 \pm 0.346}\) \\
\bottomrule
\end{tabularx}
\end{table}


The three-channel input performed slightly better than the two-window representation, while \textbf{ResNet18-3CH+BoVW (Proposed model)} gave the strongest overall internal result. This model is treated as the proposed final pipeline because it combines the main three-channel high-pass input with the ORB-BoVW handcrafted feature-fusion branch. The observed pattern is consistent with the high-pass channel and handcrafted local-feature branch adding information that was not fully captured by the CNN representation alone. The interpretation of why this pattern occurs is discussed further in Chapter~6.


\subsubsection{Uncertainty of the Proposed Model}
\label{sec:proposed_model_uncertainty}

To give an additional indication of robustness for the proposed model, Table~\ref{tab:proposed_model_ci} reports approximate 95\% confidence intervals across the five internal validation folds. The intervals are computed from the fold mean and standard deviation using a \(t\)-interval with four degrees of freedom. Since only five folds are available, the intervals should be interpreted cautiously, but they provide a compact summary of the fold-level uncertainty for the final model.

\begin{table}[H]
\centering
\caption{Approximate 95\% confidence intervals for \textbf{ResNet18-3CH+BoVW (Proposed model)} across the five internal validation folds.}
\label{tab:proposed_model_ci}
\small
\begin{tabular}{lcc}
\toprule
Metric & Mean \(\pm\) SD & Approx. 95\% CI \\
\midrule
MAE & \(\mathbf{1.244 \pm 0.297}\) & \([0.875,\ 1.613]\) \\
RMSE & \(\mathbf{1.712 \pm 0.505}\) & \([1.085,\ 2.339]\) \\
\(R^2\) & \(\mathbf{0.236 \pm 0.295}\) & \([-0.130,\ 0.602]\) \\
Spearman & \(\mathbf{0.498 \pm 0.346}\) & \([0.068,\ 0.928]\) \\
\bottomrule
\end{tabular}
\end{table}


\subsection{Backbone Comparison}
\label{sec:backbone_comparison_results}

Table~\ref{tab:backbone_comparison_results} compares ResNet18 and DenseNet121 across the same internal model settings: two-window input, three-channel input, and ORB-BoVW fusion. The ResNet18 rows repeat the corresponding results from Table~\ref{tab:internal_representation_fusion_results} so that the backbone comparison can be read directly. In this comparison, ResNet18 uses the ImageNet-based initialization described in Chapter~4, while DenseNet121 uses pretrained RadImageNet weights. All models were trained on the log-transformed target, and predictions were transformed back to the raw mucus-score scale for reporting.


\begin{table}[H]
\centering
\caption{Backbone comparison across input representation and fusion settings. Values are reported as mean \(\pm\) standard deviation across patient-level folds.}
\label{tab:backbone_comparison_results}
\small
\begin{tabular}{llcc}
\toprule
Setting & Metric & ResNet18 & DenseNet121 \\
\midrule
Two-window input & MAE & \(1.451 \pm 0.395\) & \(1.602 \pm 0.344\) \\
 & RMSE & \(1.858 \pm 0.434\) & \(2.038 \pm 0.483\) \\
 & \(R^2\) & \(0.126 \pm 0.161\) & \(-0.086 \pm 0.168\) \\
 & Spearman & \(0.363 \pm 0.239\) & \(0.208 \pm 0.239\) \\
\midrule
Three-channel input & MAE & \(1.386 \pm 0.456\) & \(1.356 \pm 0.237\) \\
 & RMSE & \(1.826 \pm 0.394\) & \(1.789 \pm 0.384\) \\
 & \(R^2\) & \(0.186 \pm 0.394\) & \(0.084 \pm 0.126\) \\
 & Spearman & \(0.387 \pm 0.256\) & \(0.247 \pm 0.189\) \\
\midrule
ORB-BoVW fusion & MAE & \(1.244 \pm 0.297\) & \(1.422 \pm 0.354\) \\
 & RMSE & \(1.712 \pm 0.505\) & \(1.850 \pm 0.566 \) \\
 & \(R^2\) & \(0.236 \pm 0.295\) & \(0.107 \pm 0.322\) \\
 & Spearman & \(0.498 \pm 0.346\) & \(0.408 \pm 0.358\) \\
\bottomrule
\end{tabular}
\end{table}

DenseNet121 gave slightly lower MAE and RMSE in the three-channel setting, but ResNet18 performed better in the two-window and fusion settings and had stronger rank correlation overall. Therefore, the results do not show a consistent advantage from replacing ResNet18 with DenseNet121.



\section{Transfer Learning Experiments}
\label{sec:transfer_learning_experiments}

The transfer-learning experiments use two source-model variants from the main pipeline. The first variant, \textbf{ResNet18-3CH} corresponds to the ResNet18 slice-based model with the main three-channel high-pass input representation and the main configuration settings described in Chapter~4. The second variant, \textbf{ResNet18-3CH+BoVW}, is the same model but with the ORB-BoVW handcrafted feature-fusion branch added. The fusion variant refers to the \textbf{ResNet18-3CH+BoVW (Proposed model)}, not to a separate standalone handcrafted model.

In the tables, zero-shot means that the source-domain checkpoint is evaluated directly on the target domain without target-domain training. Fine-tuned means that the source-domain checkpoint is used as initialization and then further trained on the target-domain training data before evaluation.

Because the MosMed target differs from the internal mucus-burden target, the transfer tables should be read as evidence about source--target adaptation between CT tasks rather than as clinical external validation for mucus plugging.


\subsection{Mucus-to-MosMed Transfer}
\label{sec:mucus_to_mosmed_transfer}

Table~\ref{tab:transfer_mosmed_target} reports transfer-learning results where the mucus dataset is used as the source domain and MosMed is used as the target domain. In this direction, the target-domain evaluation is performed on the MosMed case-level severity target.

\begin{table}[H]
\centering
\caption{Transfer results from mucus source domain to MosMed target domain.}
\label{tab:transfer_mosmed_target}
\begin{tabular}{lcccc}
\toprule
Transfer setting & MAE & RMSE & \(R^2\) & Spearman \\
\midrule
ResNet18-3CH, zero-shot & 0.903 & 1.070 & -1.62 & -0.027 \\
ResNet18-3CH, fine-tuned on MosMed & 0.460 & 0.643 & 0.04 & 0.320 \\
ResNet18-3CH+BoVW, zero-shot & 0.856 & 1.013 & -0.95 & -0.160 \\
ResNet18-3CH+BoVW, fine-tuned on MosMed & 0.369 & 0.545 & 0.29 & 0.489 \\
\bottomrule
\end{tabular}
\end{table}

The fine-tuned setting performed better than zero-shot evaluation for both model variants. The fine-tuned \textbf{ResNet18-3CH+BoVW (Proposed model)} obtained the lowest MAE and RMSE, and also the highest \(R^2\) and Spearman correlation in this transfer direction.


\subsection{MosMed-to-Mucus Transfer}
\label{sec:mosmed_to_mucus_transfer}

Table~\ref{tab:transfer_mucus_target} reports transfer-learning results where MosMed is used as the source domain and the mucus dataset is used as the target domain. In this direction, the target-domain evaluation is performed on the internal mucus-burden score.

\begin{table}[H]
\centering
\caption{Transfer results from MosMed source domain to mucus target domain.}
\label{tab:transfer_mucus_target}
\begin{tabular}{lcccc}
\toprule
Transfer setting & MAE & RMSE & \(R^2\) & Spearman \\
\midrule
ResNet18-3CH, zero-shot & 1.889 & 2.199 & 0.08 & 0.050 \\
ResNet18-3CH, fine-tuned on mucus & 1.287 & 1.576 & 0.22 & 0.139 \\
ResNet18-3CH+BoVW, zero-shot & 1.802 & 2.097 & 0.10 & 0.261 \\
ResNet18-3CH+BoVW, fine-tuned on mucus & 1.329 & 1.605 & 0.12 & 0.304 \\
\bottomrule
\end{tabular}
\end{table}

The fine-tuned setting performed better than zero-shot transfer on the mucus target for both model variants. The fine-tuned ResNet18-3CH model obtained the lowest MAE, lowest RMSE, and highest \(R^2\), while the fine-tuned \textbf{ResNet18-3CH+BoVW (Proposed model)} obtained the highest Spearman correlation.


\section{Result Summary}
\label{sec:result_summary}
The baseline results show that the representative slice-selection strategy performed better than uniform sampling, with lower MAE and RMSE and a higher Spearman correlation. This gives a simple reference point before comparing the full internal model variants.

In the main internal experiments with the final pipeline configuration, the \textbf{ResNet18-3CH+BoVW (Proposed model)} obtained the lowest MAE and RMSE and the highest \(R^2\) and Spearman correlation among the listed internal experiments. The three-channel input also performed better than the two-window input, although the strongest internal result was obtained when feature fusion was added.

The backbone comparison shows that DenseNet121 did not consistently outperform ResNet18 across the reported settings. DenseNet121 gave slightly lower MAE and RMSE in the three-channel setting, but ResNet18 performed better in the two-window and fusion settings and had stronger rank correlation overall.

In the transfer-learning experiments, the fine-tuned settings performed better than zero-shot transfer in both directions. The transfer results are reported as average performance across the five folds. In the mucus-to-MosMed direction, the fine-tuned \textbf{ResNet18-3CH+BoVW (Proposed model)} performed best across all reported metrics for the MosMed severity target. In the MosMed-to-mucus direction, the fine-tuned ResNet18-3CH model obtained the lowest MAE, lowest RMSE, and highest \(R^2\), while the fine-tuned \textbf{ResNet18-3CH+BoVW (Proposed model)} obtained the highest Spearman correlation. These findings describe transfer behaviour between different CT prediction tasks and should not be interpreted as external mucus-specific validation.

Several \(R^2\) values are close to zero or negative, especially in the baseline and zero-shot transfer settings. These values are kept in the tables because they show the difficulty of the task and the limited usefulness of direct transfer without target-domain adaptation. They also support interpreting the results as exploratory patient-level prediction results rather than definitive clinical evidence.

The deeper interpretation of these results, including the role of slice selection, high-pass input representation, feature fusion, backbone choice, and transfer direction, is discussed in Chapter~6.
